\documentclass[11pt]{article}
\usepackage[T1]{fontenc}
\usepackage{lmodern,amsmath,amssymb,amsthm,mathtools}
\usepackage[margin=1in]{geometry}
\usepackage[colorlinks=true,linkcolor=blue,citecolor=blue,urlcolor=blue]{hyperref}
\hypersetup{pdftitle={Minimax private density estimation with mixed anisotropy},pdfauthor={Kaiwen Jin}}
\newtheorem{theorem}{Theorem}
\newtheorem{lemma}[theorem]{Lemma}
\newtheorem{proposition}[theorem]{Proposition}
\newcommand{\E}{\mathbb E}
\newcommand{\R}{\mathbb R}
\newcommand{\N}{\mathbb N}
\newcommand{\one}{\mathbf 1}
\newcommand{\TV}{\operatorname{TV}}
\newcommand{\KL}{\operatorname{KL}}
\newcommand{\ind}{\mathbf 1}
\title{Minimax private density estimation with mixed anisotropy}
\author{Kaiwen Jin\\\small\href{mailto:yui.ui.siki@gmail.com}{yui.ui.siki@gmail.com}}
\date{October 13, 2026}
\begin{document}
\maketitle

\begin{abstract}
We determine the minimax density-estimation rate under sequentially interactive
local differential privacy and an anisotropic negative Sobolev loss, when the
density and discriminator smoothness vectors are proportional. The result
covers coordinates whose discriminator smoothness lies on opposite sides of
the ambient dimension. The harmonic mean of discriminator smoothness determines
the rate, with a single logarithmic factor at the critical value. A weighted
randomized response for individual Fourier coefficients gives the upper bound.
The noncritical lower bounds follow from Albert, Chevallier, Laurent and Sacko;
at the critical value, an orthogonal information bound for the full private
history yields the matching logarithmic lower bound.
\end{abstract}

\section{The model and the rate}

Local differential privacy requires each individual to release a randomized
observation whose law changes by at most a prescribed likelihood ratio when
that individual's data change. Density estimation under this constraint can
be assessed through smooth test functions: the loss measures how well the
estimated distribution reproduces their expectations. Different coordinates
may require different degrees of smoothness.

Albert, Chevallier, Laurent and Sacko~\cite{ACLS} study this problem with Fourier
Sobolev density and discriminator classes. Their Theorem~10 gives an
anisotropic lower bound when the two smoothness vectors are proportional.
Their Theorem~11 gives upper bounds when all discriminator smoothness
coordinates lie below the dimension, all lie above it, or all equal it.
The paragraph following that theorem, on p.~15, leaves the mixed case open.
Here we retain the sum of the anisotropic Fourier weights and sample
coefficients with probabilities adapted to that sum. This supplies an upper
bound in the mixed case and leads to a precise critical logarithm.
The noncritical lower theorem is used with its original hypotheses and credit;
the critical lower argument is given below.

Fix an integer $d\ge2$, $R^2>d$, $A>0$, and vectors
$\boldsymbol\beta,\boldsymbol\delta\in\N_{>0}^d$ such that
$\beta_m/\delta_m=c>0$ is independent of $m$. Assume that some
$\delta_m<d$ and some $\delta_k>d$. Other coordinates may equal $d$.
On $[0,1]$, use the real Fourier basis
\[
 \phi_1(t)=1,\qquad
 \phi_{2j}(t)=\sqrt2\cos(2\pi jt),\qquad
 \phi_{2j+1}(t)=\sqrt2\sin(2\pi jt),\quad j\ge1.
\]
Its tensor products are $\phi_{\mathbf j}(x)=\prod_{m=1}^d\phi_{j_m}(x_m)$,
with $\mathbf j\in\N_{>0}^d$. All integrals below use Lebesgue measure on
$[0,1]^d$. Write $\theta_{\mathbf j}(h)=\int h\phi_{\mathbf j}$ and, for a
positive smoothness vector $\mathbf s$, define
\[
 W^{\mathbf s}(r)=\left\{h\in L^2([0,1]^d):
 \sum_{\mathbf j}\left(\sum_{m=1}^d j_m^{2s_m}\right)
 |\theta_{\mathbf j}(h)|^2\le r^2\right\}.
\]
Let $\mathcal F_{\boldsymbol\beta,R}$ be the probability densities in
$W^{\boldsymbol\beta}(R)$, and set
\[
 d_{\boldsymbol\delta}(h,f)=
 \sup_{g\in W^{\boldsymbol\delta}(1)}\left|\int(h-f)g\right|.
\]
The data $X_1,\ldots,X_n$ are independent with common density $f$.
A sequentially interactive mechanism releases $Z_i$ through a probability
kernel $Q_i(\cdot\mid X_i,Z_{<i})$. Its $\alpha$-privacy requirement is
\begin{equation}\label{eq:privacy}
 Q_i(B\mid x,z_{<i})\le e^\alpha Q_i(B\mid x',z_{<i})
\end{equation}
for every measurable output event $B$, every $x,x'$, and every complete
history $z_{<i}$. Denote the class of such mechanisms by $\mathcal Q_\alpha$.
The risk of interest is
\begin{equation}\label{eq:risk}
 \mathcal R_{n,\alpha}(\boldsymbol\beta,\boldsymbol\delta)
 =\inf_{Q\in\mathcal Q_\alpha}\inf_{\widehat f}
  \sup_{f\in\mathcal F_{\boldsymbol\beta,R}}
  \E_{f,Q}d_{\boldsymbol\delta}(\widehat f,f),
\end{equation}
where $\widehat f$ is a measurable $L^2$-valued function of the released
data. The loss in~\eqref{eq:risk} is an expected norm, without squaring.

\begin{theorem}\label{thm:main}
Under the fixed parameter assumptions above, put
\[
 T=n\alpha^2,\qquad q=\sum_{m=1}^d\frac1{\delta_m},\qquad
 c=\frac{\beta_m}{\delta_m}.
\]
There are constants $C_-,C_+>0$, depending only on
$d,R,A,\boldsymbol\beta,\boldsymbol\delta$, such that for every integer
$n\ge1$ and every $0<\alpha\le A$ with $T\ge1$,
\begin{equation}\label{eq:mainrate}
 C_-r(T)\le\mathcal R_{n,\alpha}(\boldsymbol\beta,\boldsymbol\delta)
 \le C_+r(T),\qquad
 r(T)=\begin{cases}
 T^{-(c+1)/(2(c+q))},&q>1,\\[2pt]
 T^{-1/2}\log(eT),&q=1,\\[2pt]
 T^{-1/2},&q<1.
 \end{cases}
\end{equation}
\end{theorem}

Thus the three cases correspond to the harmonic mean
$\bar\delta=d/q$ lying below, at, or above $d$. The power exponent can also
be written $(\bar\beta+\bar\delta)/(2\bar\beta+2d)$, where
$\bar\beta^{-1}=d^{-1}\sum_m\beta_m^{-1}$.
The theorem concerns known, fixed smoothness parameters and pure local
privacy. No adaptation to unknown smoothness is asserted.

\section{Fourier weights and a private estimator}

We suppress boldface on a Fourier multi-index $j$. Define
\[
 u_j=\sum_m j_m^{2\beta_m},\qquad v_j=\sum_m j_m^{2\delta_m}.
\]
The constant index is $\one=(1,\ldots,1)$. A density has
$\theta_{\one}=1$ and $u_{\one}=v_{\one}=d$.
For every $h,f\in L^2$ the exact dual norm is
\begin{equation}\label{eq:dual}
 d_{\boldsymbol\delta}(h,f)^2
 =\sum_j\frac{|\theta_j(h-f)|^2}{v_j}.
\end{equation}
Weighted Cauchy--Schwarz gives one inequality. For a finite set $E$, the
Fourier polynomial with coefficients proportional to
$\theta_j(h-f)/v_j$, $j\in E$, attains the square root of the corresponding
partial sum after normalization in $W^{\boldsymbol\delta}(1)$.
Increasing finite sets proves the reverse inequality. The full sum is finite
because $v_j\ge d$. Also, comparison with the largest summand gives
\begin{equation}\label{eq:weights}
 d^{-c}v_j^c\le u_j\le d v_j^c.
\end{equation}

The coefficient sampling probabilities will depend on the following sum.
For $t>0$, let
\[
 E_t=\{j\ne\one:v_j\le t\},\qquad
 N(t)=|E_t|,\qquad S(t)=\sum_{j\in E_t}v_j^{-1/2}.
\]
Throughout the proof, comparison constants depend only on the fixed
parameters, never on $n$, $\alpha$, or the chosen mechanism.

\begin{lemma}\label{lem:count}
For all sufficiently large $t$,
\[
 N(t)\asymp t^{q/2},\qquad
 S(t)\asymp\begin{cases}
 t^{(q-1)/2},&q>1,\\
 \log t,&q=1,\\
 1,&q<1.
 \end{cases}
\]
\end{lemma}
\begin{proof}
The inequalities $j_m\le t^{1/(2\delta_m)}$ give the upper count.
For the lower count, every point in the box
\[
 1\le j_m\le\left\lfloor(t/d)^{1/(2\delta_m)}\right\rfloor,
 \qquad 1\le m\le d,
\]
has $v_j\le t$. For large $t$, each floor is at least half its argument;
removing the constant index leaves a count comparable to $t^{q/2}$.

Dyadic shells and the upper count bound the contribution at weight scale $s$
by a constant times $s^{(q-1)/2}$. Summing these bounds gives the three upper
estimates. To obtain lower estimates, let $a_0t^{q/2}\le N(t)\le a_1t^{q/2}$
for large $t$ and choose a fixed $H>1$ with
$a_1H^{-q/2}\le a_0/2$. Then
\[
 N(t)-N(t/H)\ge (a_0/2)t^{q/2}
\]
once both arguments are large. The shell $(t/H,t]$ therefore contributes
at least $(a_0/2)t^{(q-1)/2}$ to $S(t)$. For $q>1$ use the top shell.
For $q=1$ use the disjoint shells down to a fixed threshold; their number is
comparable to $\log t$. For $q<1$, one fixed nonconstant index gives the
required positive lower bound.
\end{proof}

Choose a fixed $t_0$ large enough for Lemma~\ref{lem:count} and for
$E_{t_0}\ne\varnothing$. On $E_t$, $t\ge t_0$, set
\begin{equation}\label{eq:sampling}
 p_j=\frac{v_j^{-1/2}}{S(t)},\qquad
 B_0=2^{d/2},\qquad \kappa=\frac{e^\alpha-1}{e^\alpha+1}.
\end{equation}
An individual draws $J$ with law $(p_j)$, independently of $X$, and releases
$(J,Y)$, where $Y\in\{-1,1\}$ has conditional law
\begin{equation}\label{eq:channel}
 \Pr(Y=y\mid X=x,J=j)=\frac12
 \left(1+\frac{y\kappa\phi_j(x)}{B_0}\right).
\end{equation}
Use independent randomization for the $n$ individuals and define
\begin{equation}\label{eq:estimator}
 \widehat\theta_j=\frac{B_0}{n\kappa p_j}
 \sum_{i=1}^n\ind_{\{J_i=j\}}Y_i,\qquad
 \widehat f=1+\sum_{j\in E_t}\widehat\theta_j\phi_j.
\end{equation}
This estimator is a measurable finite Fourier polynomial.

\begin{proposition}\label{prop:upper}
The channel~\eqref{eq:channel} belongs to $\mathcal Q_\alpha$, and
\begin{equation}\label{eq:upper}
 \sup_{f\in\mathcal F_{\boldsymbol\beta,R}}
 \E_{f,Q}d_{\boldsymbol\delta}(\widehat f,f)
 \le C\left(t^{-(c+1)/2}+\frac{S(t)}{\sqrt T}\right).
\end{equation}
\end{proposition}
\begin{proof}
Since $|\phi_j(x)|\le B_0$, the ratio of any output-atom probabilities for
two input points is at most
$(1+\kappa)/(1-\kappa)=e^\alpha$. Summing over atoms proves privacy for
all events. The channel is independent of the preceding history, so this
inequality holds for every history in~\eqref{eq:privacy}.

Each coefficient in~\eqref{eq:estimator} is unbiased. The squared single
observation contribution has expectation $B_0^2/(\kappa^2p_j)$, hence
\[
 \operatorname{Var}(\widehat\theta_j)\le\frac{B_0^2}{n\kappa^2p_j},
 \qquad
 \sum_{j\in E_t}\frac{\operatorname{Var}(\widehat\theta_j)}{v_j}
 \le\frac{B_0^2 S(t)^2}{n\kappa^2}.
\]
The diagonal sum in~\eqref{eq:dual} requires no independence between the
estimated coordinates. The omitted-coefficient bias, by~\eqref{eq:weights}, is
\[
 \sum_{v_j>t}\frac{\theta_j(f)^2}{v_j}
 \le d^c t^{-(c+1)}\sum_j u_j\theta_j(f)^2
 \le d^c R^2t^{-(c+1)}.
\]
Finally, $\kappa\ge\alpha/(e^A+1)$ because
$e^\alpha-1\ge\alpha$ and $e^\alpha+1\le e^A+1$.
Apply Cauchy--Schwarz to the expectation of the norm to obtain
\eqref{eq:upper}.
\end{proof}

For $q>1$, take $t=t_0T^{1/(c+q)}$. For $q\le1$, take
$t=t_0T^{1/(c+1)}$. Lemma~\ref{lem:count} and
Proposition~\ref{prop:upper} give the upper bound in~\eqref{eq:mainrate},
including $T=1$.

\section{An information bound with the full private history}

The critical lower bound needs one information estimate shared by all the
coefficients. Bounding each neighboring pair separately without this shared
sum would miss the logarithm. We prove the estimate on joint transcript
spaces, which also handles arbitrary measurable output spaces.

\begin{lemma}\label{lem:trace}
Let $(\phi_j)_{j\in E}$ be a finite mean-zero orthonormal family on the
Lebesgue probability space $[0,1]^d$. Suppose $a_j>0$ and every density
\[
 f_\varepsilon=1+\sum_{j\in E}\varepsilon_j a_j\phi_j,
 \qquad \varepsilon\in\{-1,1\}^{E},
\]
lies between $1/2$ and $3/2$ almost everywhere. Fix any
$Q\in\mathcal Q_\alpha$, and let $M_\varepsilon^n$ be its full transcript
law under $f_\varepsilon$. If $\varepsilon^{(j)}$ denotes the vector with
sign $j$ reversed, there are nonnegative numbers $\bar I_j$ such that
\begin{equation}\label{eq:trace-test}
 \sum_{j\in E}\bar I_j\le n(e^\alpha-1)^2,\qquad
 \tau_j:=2^{-|E|}\sum_\varepsilon
 \TV(M_\varepsilon^n,M_{\varepsilon^{(j)}}^n)
 \le2a_j\sqrt{\bar I_j}.
\end{equation}
Here $\TV(P,Q)=\sup_B|P(B)-Q(B)|$.
\end{lemma}
\begin{proof}
Write $M_0^i$ for the law of the first $i$ releases under the uniform density,
and let $h=z_{<i}$ denote the full history. Its uniform-input next-step kernel
is $K_{0,i}(\cdot\mid h)=\int Q_i(\cdot\mid x,h)\,dx$.
Integrating~\eqref{eq:privacy} over a uniform second input gives, for every
$x,h$,
\begin{equation}\label{eq:kernel-dom}
 e^{-\alpha}K_{0,i}(\cdot\mid h)\le Q_i(\cdot\mid x,h)
 \le e^\alpha K_{0,i}(\cdot\mid h).
\end{equation}

Introduce two finite measures on the joint space of $(x,h,z)$:
\[
 \Gamma_i(dx,dh,dz)=dx\,M_0^{i-1}(dh)Q_i(dz\mid x,h),\qquad
 \Lambda_i(dx,dh,dz)=dx\,M_0^i(dh,dz).
\]
Applying~\eqref{eq:kernel-dom} to measurable sections and integrating gives
$e^{-\alpha}\Lambda_i\le\Gamma_i\le e^\alpha\Lambda_i$.
The Radon--Nikodym theorem and Fubini's theorem for finite measures
\cite[232F, 252C, 252H]{Fremlin} give a jointly measurable density
$r_i=d\Gamma_i/d\Lambda_i$ satisfying
\[
 e^{-\alpha}\le r_i\le e^\alpha,\qquad
 \int r_i(x,h,z)\,dx=1
\]
for $M_0^i$-almost every $(h,z)$. The last equality follows because the
projection of $\Gamma_i$ onto $(h,z)$ is $M_0^i$.
Set
\[
 s_{ij}(h,z)=\int\phi_j(x)r_i(x,h,z)\,dx,\qquad
 m_{i,\varepsilon}(h,z)=1+\sum_j\varepsilon_j a_j s_{ij}(h,z).
\]
Then $m_{i,\varepsilon}=\int f_\varepsilon r_i\,dx\in[1/2,3/2]$.
Expanding the integral of the nonnegative square
$(r_i-1-\sum_j s_{ij}\phi_j)^2$, using the mean-zero orthonormal family,
yields the pointwise joint bound
\begin{equation}\label{eq:orth-trace}
 \sum_j s_{ij}(h,z)^2\le\int(r_i-1)^2dx\le(e^\alpha-1)^2.
\end{equation}
For each bounded measurable function $F$ of the preceding history,
\begin{equation}\label{eq:score-zero}
 \int F(h)s_{ij}(h,z)\,dM_0^i(h,z)=0.
\end{equation}
Indeed, use $\Gamma_i$ for this integral and integrate $z$ first; the
remaining integral of $\phi_j$ over $x$ vanishes. The identity extends to
integrable history weights by truncation, since $s_{ij}$ is bounded.

Let $L_\varepsilon^i=dM_\varepsilon^i/dM_0^i$.
The fresh $X_i$ is independent of the preceding released history under each
fixed density. Induction therefore gives
\begin{equation}\label{eq:sequential-density}
 L_\varepsilon^i(h,z)
 =L_\varepsilon^{i-1}(h)m_{i,\varepsilon}(h,z),\qquad
 L_\varepsilon^0=1.
\end{equation}
To check the identity, test its right-hand side against a bounded measurable
function of $(h,z)$, replace $m_{i,\varepsilon}$ by
$\int f_\varepsilon r_i dx$, and use $\Gamma_i$ and the preceding
history density. This gives exactly the next-step law under $f_\varepsilon$.
All $L_\varepsilon^i$ are positive. This argument does not require measurable
versions of conditional Radon--Nikodym kernels on the output space.

For neighboring signs,
$m_{i,\varepsilon}-m_{i,\varepsilon^{(j)}}
=2\varepsilon_j a_j s_{ij}$. By~\eqref{eq:sequential-density}, the logarithm
of the full likelihood ratio is the sum of the corresponding logarithms of
the $m$'s. For positive $a,b$,
\[
 a\log(a/b)\le\frac{(a-b)^2}{b}+(a-b).
\]
The last term has zero integral against the preceding history density by
\eqref{eq:score-zero}. The denominator $m_{i,\varepsilon^{(j)}}$ is at
least $1/2$, so
\begin{equation}\label{eq:kl-bound}
 \KL(M_\varepsilon^n,M_{\varepsilon^{(j)}}^n)
 \le8a_j^2 I_j(\varepsilon),\qquad
 I_j(\varepsilon)=\sum_{i=1}^n\int
 L_\varepsilon^{i-1}(h)s_{ij}(h,z)^2\,dM_0^i(h,z).
\end{equation}
All logarithms are bounded at each step, and all these integrals are finite.
Since $\int L_\varepsilon^{i-1}(h)dM_0^i=1$, summing
\eqref{eq:orth-trace} gives
\[
 \sum_jI_j(\varepsilon)\le n(e^\alpha-1)^2.
\]
Define $\bar I_j$ as the uniform average of $I_j(\varepsilon)$.

In our normalization, $\TV(P,Q)^2\le\KL(P,Q)/2$.
Here this inequality can be seen directly by taking
$B=\{dP/dQ\ge1\}$ and writing $p=P(B)$, $q_0=Q(B)$.
The convex tangent inequality for $t\log t$ on $B$ and its complement gives
\[
 \KL(P,Q)\ge p\log\frac p{q_0}
 +(1-p)\log\frac{1-p}{1-q_0}\ge2(p-q_0)^2.
\]
For the second inequality the expression as a function of $p$ has value and
first derivative zero at $q_0$, and second derivative
$1/[p(1-p)]\ge4$. Zero-mass endpoints are handled by continuity or infinite
relative entropy. Also $p-q_0=\TV(P,Q)$.
Apply this inequality to~\eqref{eq:kl-bound} and use Cauchy--Schwarz when
averaging over signs. This proves~\eqref{eq:trace-test}.
\end{proof}

The joint measures in this proof keep every preceding release in the
history. The private releases need not be independent. The argument refines
the use of privacy contraction in local-private minimax bounds
\cite{DJW}, by retaining a sum over orthogonal coefficients at each step.

\section{The critical lower bound and completion of the proof}

At $q=1$, Lemma~\ref{lem:count} supplies a logarithm in the weighted spectral
sum. We put a finite family of densities on those coefficients, with
amplitudes that make the testing bound spend the same shared information sum.
The cutoff can grow slowly: the lower bound needs a number of scales
proportional to $\log T$, rather than the largest admissible cutoff.

Assume $q=1$ and set
\[
 \rho=\frac1{4(c+1)},\qquad t=T^\rho,\qquad E=E_t,\qquad S=S(t).
\]
For $T$ sufficiently large, define positive amplitudes by
\begin{equation}\label{eq:amplitudes}
 a_j^2=\frac{\eta^2 S\sqrt{v_j}}{T},\qquad j\in E,
\end{equation}
where $\eta>0$ is fixed below. Lemma~\ref{lem:count} gives
$N(t)\lesssim t^{1/2}$ and $S\asymp\log t\asymp\log T$.
Consequently,
\begin{align}
 B_0\sum_{j\in E}a_j&\lesssim
 \eta\sqrt{S/T}\,t^{3/4}\longrightarrow0,\label{eq:positive}\\
 \sum_{j\in E}u_ja_j^2&\lesssim
 \eta^2(S/T)t^{c+1}
 =\eta^2 S T^{-3/4}\longrightarrow0.\label{eq:smooth}
\end{align}
The first convergence uses $\rho<1/4$, and the second uses
\eqref{eq:weights}. Since $R^2>d$, there is a fixed $T_*$ such that for
every $T\ge T_*$ and every sign vector the Fourier polynomial
\[
 f_\varepsilon=1+\sum_{j\in E}\varepsilon_ja_j\phi_j
\]
lies between $1/2$ and $3/2$ and has Sobolev weight
$d+\sum_j u_ja_j^2\le R^2$. Each nonconstant basis function has mean zero,
so all these polynomials are probability densities in the original class.
The positivity is uniform over every sign vector.

Fix any mechanism $Q\in\mathcal Q_\alpha$ and any measurable $L^2$-valued
estimator. Decode its $j$th Fourier coefficient as $+1$ when it is
nonnegative and as $-1$ otherwise. This is measurable, and a wrong decoded
sign causes coefficient error at least $a_j$. Write
\[
 w_j=v_j^{-1},\qquad b_j=w_ja_j^2,\qquad V=\sum_{j\in E}b_j.
\]
By~\eqref{eq:dual}, pointwise,
\begin{equation}\label{eq:norm-hamming}
 d_{\boldsymbol\delta}(\widehat f,f_\varepsilon)
 \ge\sqrt{\sum_j b_j\ind_{\{\widehat\varepsilon_j\ne\varepsilon_j\}}}
 \ge\frac1{\sqrt V}\sum_jb_j\ind_{\{\widehat\varepsilon_j\ne\varepsilon_j\}}.
\end{equation}
The second inequality uses $\sqrt x\ge x/\sqrt V$ for $0\le x\le V$.
For a neighboring pair of transcript laws, the sum of the two decision
error probabilities is at least $1-\TV$.
Thus the sign-average risk is at least
\begin{equation}\label{eq:average-risk}
 \frac{V-\sum_j b_j\tau_j}{2\sqrt V},
\end{equation}
where $\tau_j$ is from Lemma~\ref{lem:trace}. This is a lower bound for the
expected norm itself.

For the amplitudes~\eqref{eq:amplitudes}, direct summation gives
\[
 V=\frac{\eta^2 S^2}{T},\qquad
 \sum_j b_j^2a_j^2=\frac{\eta^6 S^4}{T^3}.
\]
Moreover $e^\alpha-1\le e^A\alpha$. Weighted Cauchy--Schwarz and
\eqref{eq:trace-test} therefore yield
\[
 \sum_j b_j\tau_j
 \le2\left(\sum_j b_j^2a_j^2\right)^{1/2}
       \left(\sum_j\bar I_j\right)^{1/2}
 \le\frac{2e^A\eta^3 S^2}{T}.
\]
Choose $0<\eta\le1/(4e^A)$. The last quantity is at most $V/2$.
Equation~\eqref{eq:average-risk} now gives a lower bound
$\eta S/(4\sqrt T)\gtrsim\log T/\sqrt T$ for every mechanism and estimator,
uniformly for $T\ge T_*$. The finite density family and all constants were
chosen independently of the mechanism.

\begin{proof}[Proof of Theorem~\ref{thm:main}]
The upper bound was established after Proposition~\ref{prop:upper}.
For the remaining lower estimates, Theorem~10 of~\cite[p.~14]{ACLS} states,
under the same positive integer smoothness, proportionality, radius and
privacy assumptions,
\[
 \mathcal R_{n,\alpha}(\boldsymbol\beta,\boldsymbol\delta)
 \ge C\max\left\{
 T^{-(\bar\beta+\bar\delta)/(2\bar\beta+2d)},T^{-1/2}\right\}.
\]
Its mechanism class uses the full-history privacy condition
\eqref{eq:privacy}; its minimax loss is the expected unsquared adversarial
distance. The absolute value in our definition gives the same distance,
since the discriminator ball is symmetric. The real Fourier indexing and
sum weights are also the same.
As $\bar\delta=d/q$ and $\bar\beta=cd/q$, the first exponent is
$(c+1)/(2(c+q))$. This proves the required lower bound for $q>1$ and $q<1$.
When $q=1$, the preceding construction proves the logarithmic lower bound
for $T\ge T_*$. For $1\le T<T_*$, the cited bound is at least
$C T^{-1/2}$, and $\log(eT)\le\log(eT_*)$; decreasing the fixed lower
constant covers this interval as well. Both lower arguments hold for every
$Q$ and every measurable $L^2$-valued estimator before taking the infima in
\eqref{eq:risk}. This proves~\eqref{eq:mainrate} throughout the stated
parameter range.
\end{proof}

\paragraph{AI-use statement.}
OpenAI Codex assisted with mathematical exploration, proof drafting,
reference checks and manuscript preparation. Independent mathematical review
is requested.

\begin{thebibliography}{9}
\bibitem{ACLS}
M. Albert, J. Chevallier, B. Laurent and O. Sacko.
\newblock Minimax density estimation in the adversarial framework under local
differential privacy.
\newblock \emph{Journal of Machine Learning Research} \textbf{27}(69)
(2026), 1--43.
\newblock \url{https://www.jmlr.org/papers/v27/24-0494.html}.

\bibitem{DJW}
J. C. Duchi, M. I. Jordan and M. J. Wainwright.
\newblock Minimax optimal procedures for locally private estimation.
\newblock \emph{Journal of the American Statistical Association}
\textbf{113}(521) (2018), 182--201.
\newblock \href{https://doi.org/10.1080/01621459.2017.1389735}{doi:10.1080/01621459.2017.1389735}.
\newblock Author manuscript: \href{https://arxiv.org/abs/1604.02390v2}{arXiv:1604.02390v2},
Theorem~1 and Section~5.1.

\bibitem{Fremlin}
D. H. Fremlin.
\newblock \emph{Measure Theory}, volume 2, author-hosted abridged text,
chapters 23 and 25; Sections 232F, 252C and 252H.
\newblock \url{https://www1.essex.ac.uk/maths/people/fremlin/chap23.ro.pdf};
\url{https://www1.essex.ac.uk/maths/people/fremlin/chap25.ro.pdf}.
\end{thebibliography}
\end{document}
